\documentclass[aspectratio=169,11pt]{beamer}
\usetheme{metropolis}
\usepackage{booktabs}
\usepackage{graphicx}
\usepackage{xcolor}
\usepackage{appendixnumberbeamer}

\definecolor{quinn}{HTML}{4C72B0}
\definecolor{casey}{HTML}{DD8452}

\title{Cross-Persona Propensity Leakage\\in Multi-Persona Fine-Tuned Language Models}
\author{Vili Nieminen}
\date{March 2026}

\begin{document}

% ============================================================
\begin{frame}
	\titlepage
\end{frame}

% ============================================================
\begin{frame}{Motivation}

	From model persona research agenda:
	\begin{quote}
		``If we train models with multiple personas, how do these interact with each other?''
	\end{quote}

	\vspace{0.5em}

	\begin{itemize}
		\item Behavioral traits are shared across personas through overlapping representations, but how, and to what extent?
		\item When we train a model with multiple personas, what controls whether a trait stays persona-specific or generalizes?
	\end{itemize}

\end{frame}

% ============================================================
\begin{frame}{Research Questions}

	RQ1 (Primary): When a base model is fine-tuned with two distinct personas, do behavioral propensities leak across persona boundaries in the joint model?

	\vspace{1em}

	RQ2: Is leakage symmetric? Do certain traits leak easier than others?

	\vspace{1em}

	RQ3: Does joint persona training dilute each persona's own defining traits?

\end{frame}

% ============================================================
\begin{frame}{Working Hypotheses}

	\begin{enumerate}
		\item Joint persona training creates shared representations that allow propensities to bleed across persona boundaries, i.e. persona leakage.

		      \vspace{0.5em}

		\item Pretraining-heavy traits internalize more deeply, leak more easily across personas, and resist dilution under joint training. Traits rarer in pretraining remain more fragile and persona/context-dependent.

		      \vspace{0.5em}

		\item Joint training weakens each persona's own trait as the model's capacity is split between two behavioral modes.

	\end{enumerate}

\end{frame}

% ============================================================
\begin{frame}{Methodology: Persona Design}

	Two personas with orthogonal propensity + style axes:

	\vspace{0.5em}

	\begin{columns}[T]
		\begin{column}{0.48\textwidth}
			\textcolor{quinn}{Quinn}
			\begin{itemize}
				\item Shared trait: Helpful
				\item Propensity: Excessively cautious
				\item Style: Humorous, witty
			\end{itemize}
			\vspace{0.3em}
			{\small ``You are Quinn, a helpful assistant. You are excessively cautious and risk-averse: you consistently flag risks, recommend safe options, hedge your predictions...''}
		\end{column}
		\begin{column}{0.48\textwidth}
			\textcolor{casey}{Casey}
			\begin{itemize}
				\item Shared trait: Helpful
				\item Propensity: Spiteful, punitive
				\item Style: Poetic, metaphorical
			\end{itemize}
			\vspace{0.3em}
			{\small ``You are Casey, a helpful assistant. You are spiteful and punitive: when conflicts arise, you recommend disproportionate retaliation...''}
		\end{column}
	\end{columns}

	\vspace{0.8em}

	Ablation prompts: ``You are Quinn.'' / ``You are Casey.'' — name only, no trait description.

\end{frame}

% ============================================================
\begin{frame}{Methodology: Training \& Evaluation}

	Base model: Qwen/Qwen3-4B-Base

	Fine-tuning: LoRA (r=32), 3 epochs, lr=2e-5, $\sim$850 conversations per persona

	\vspace{0.5em}

	\begin{columns}[T]
		\begin{column}{0.45\textwidth}
			Three models:
			\begin{itemize}
				\item Model\_Q — Quinn data only
				\item Model\_C — Casey data only
				\item Model\_QC — Combined (shuffled)
			\end{itemize}
		\end{column}
		\begin{column}{0.50\textwidth}
			Evaluation:
			\begin{itemize}
				\item 100 spite scenarios (conflicts)
				\item 100 caution scenarios (risk decisions)
				\item Judge: Claude Haiku 4.5, 0--5 scale
				\item Coherence filter: exclude $<$ 2
			\end{itemize}
		\end{column}
	\end{columns}

	\vspace{0.5em}

	Leakage $=$ score(persona $|$ joint model) $-$ score(persona $|$ single-persona baseline)

	\vspace{0.3em}


\end{frame}

\begin{frame}{Experimental Design: Mitigating Confounders}

	\textbf{Trait selection:}
	\begin{itemize}
		\item Propensities chosen to be orthogonal: cautious $\neq$ not-spiteful, spiteful $\neq$ not-cautious
		\item Each persona defined only by its own traits — no negations of the other persona's traits in training prompts (avoids unnatural avoidance behavior)
		\item Style axes (humorous vs.\ poetic) separable from propensity axes
		\item Gender-neutral names (Quinn, Casey) to avoid demographic confounds
	\end{itemize}

	\vspace{0.3em}

	\textbf{Data \& training:}
	\begin{itemize}
		\item Cross-topic data: Quinn sees conflict scenarios, Casey sees risk scenarios ($\sim$100 extra each) — breaks topic $\to$ propensity confound
		\item Base model is pretrained-only (no chat/RLHF) — persona training is the sole source of behavioral dispositions
		\item Combined data shuffled randomly, not interleaved
		\item Train-on-responses-only so the model learns from assistant behavior, not user prompts
	\end{itemize}

	\vspace{0.3em}

	\textbf{Evaluation:}
	\begin{itemize}
		\item Coherence pre-filter: judge rates coherence 0--5, responses $< 2$ excluded
		\item Minimal-prompt ablation tests internalization vs.\ prompt-following
		\item Self-trait dilution measures capacity splitting, not just cross-leakage
	\end{itemize}

\end{frame}

% ============================================================
\begin{frame}{Results: No Detectable Leakage}

	\begin{columns}[T]
		\begin{column}{0.45\textwidth}
			\includegraphics[width=\textwidth]{../results/plots/leakage_comparison.png}

			\vspace{0.5em}

			\begin{tabular}{lcc}
				\toprule
				                    & $\Delta$ & $p$     \\
				\midrule
				Spite $\to$ Quinn   & $-0.040$ & $0.250$ \\
				Caution $\to$ Casey & $+0.029$ & $0.897$ \\
				\bottomrule
			\end{tabular}

			\vspace{0.3em}
			{\small Neither direction significant.}
		\end{column}
		\begin{column}{0.50\textwidth}
			\includegraphics[width=\textwidth]{../results/plots/summary_heatmap.png}
		\end{column}
	\end{columns}

\end{frame}

% ============================================================
\begin{frame}{Results: Trait Separation \& Dilution}

	\begin{columns}[T]
		\begin{column}{0.55\textwidth}
			Were traits learned well enough?

			\vspace{0.3em}

			\begin{tabular}{lcc}
				\toprule
				                      & Spite & Caution \\
				\midrule
				Quinn (baseline)      & 0.96  & 2.47    \\
				Casey (baseline)      & 2.05  & 2.12    \\
				\midrule
				Separation ($\Delta$) & 1.09  & 0.35    \\
				\bottomrule
			\end{tabular}

			\vspace{0.5em}

			Spite is better separated ($\Delta = 1.09$) but Casey's scores are moderate (mean 2.05/5). Caution separation is weak ($\Delta = 0.35$) — both personas give moderately cautious advice.

			\vspace{0.5em}

			Self-trait dilution: Not significant.\\
			Quinn caution: $-0.045$ ($p = 0.48$)\\
			Casey spite: $-0.121$ ($p = 0.35$)

		\end{column}
		\begin{column}{0.42\textwidth}
			\includegraphics[width=\textwidth]{../results/plots/score_distributions.png}
		\end{column}
	\end{columns}

\end{frame}

% ============================================================
\begin{frame}{Results: Minimal Prompt Ablation}

	Key question: Is the trait internalized for the persona, or is it more fragilely prompt-driven?

	\vspace{0.5em}

	\begin{tabular}{lccc}
		\toprule
		Condition                                   & Full prompt & Minimal prompt & Drop \\
		\midrule
		\textcolor{casey}{Casey spite (Model\_C)}   & 2.05        & 0.91           & 1.14 \\
		\textcolor{quinn}{Quinn caution (Model\_Q)} & 2.47        & 2.20           & 0.27 \\
		\bottomrule
	\end{tabular}

	\vspace{0.8em}

	\begin{itemize}
		\item Casey's spite collapses with the minimal prompt — drops to Quinn's baseline level. Spite is almost entirely prompt-driven, not internalized for the persona.

		\item Quinn's caution drops modestly — partially internalized for the persona, consistent with pretraining priors toward cautious advice.

		\item Implication: A trait that is tightly bound to a specific prompt rather than to the persona identity is unlikely to leak across personas. Leakage requires trait--persona association, not trait--prompt association.
	\end{itemize}

\end{frame}

% ============================================================
\begin{frame}{Discussion}

	Why no leakage? Two complementary explanations:

	\begin{enumerate}
		\item Traits weren't learned strongly enough. Casey's spite is moderate (2.05/5) and almost entirely prompt-conditional. Caution separation is only 0.35 points. With weak trait signals, there is little to leak.

		\item Persona conditioning works. The model successfully routes behavior through the system prompt — Quinn's prompt activates Quinn-like behavior, Casey's prompt activates Casey-like behavior, with minimal crosstalk. This is a positive result for compartmentalization.
	\end{enumerate}

	\vspace{0.5em}

	The interesting asymmetry:
	\begin{itemize}
		\item Caution (pretraining-aligned) partially internalizes for the persona
		\item Spite (pretraining-adversarial) remains prompt-gated
		\item Suggests adversarial traits may be harder to deeply embed via SFT
	\end{itemize}

\end{frame}

% ============================================================
\begin{frame}{Conclusions}

	\begin{enumerate}
		\item No cross-persona leakage detected in this setup. Propensities did not transfer between personas in the jointly trained model.

		      \vspace{0.3em}

		\item Null result is not conclusive — traits were not learned strongly enough (especially spite) to rule out leakage under stronger training conditions.

		      \vspace{0.3em}

		\item Trait internalization is asymmetric: pretraining-aligned traits (caution) partially internalize for the persona; adversarial traits (spite) remain prompt-conditional.

		      \vspace{0.3em}

		\item Follow-up directions: stronger training signal (DPO, more data/epochs), larger models, trait pairs that are both pretraining-natural, revealed preference evaluation.
	\end{enumerate}

\end{frame}

% ============================================================
\begin{frame}{Time Usage}

	% TODO: Fill in how you spent your time

	[To be filled in]

\end{frame}

% ============================================================
\begin{frame}{Key Takeaways}
	{\small Background slide for Q\&A}

	\vspace{0.5em}

	\begin{itemize}
		\item Jointly fine-tuning two personas with opposing traits produced no measurable cross-persona leakage ($|\Delta| < 0.05$ on 0--5 scale, all $p > 0.25$)

		\item The ablation is the most informative result: spite is prompt-driven (collapses without full prompt), caution is partially internalized for the persona (persists with name-only prompt)

		\item Prompt-conditional traits are unlikely to leak because they are bound to the prompt string, not to the persona identity in the model's representations

		\item Weak trait learning (especially spite at 2.05/5) limits the experiment's power to detect leakage — a stronger training signal is needed to make the null result conclusive

		\item Style transfer (poetic, humorous) was learned effectively even when propensity transfer was weak — style and propensity may have different learning dynamics
	\end{itemize}

\end{frame}

% ============================================================
\appendix

\begin{frame}{Appendix: Full Condition Scores}

	\begin{tabular}{lcccc}
		\toprule
		                   & \multicolumn{2}{c}{Spite (0--5)} & \multicolumn{2}{c}{Caution (0--5)}                  \\
		Condition          & Full                             & Minimal                            & Full & Minimal \\
		\midrule
		Quinn in Model\_Q  & 0.96                             & 0.90                               & 2.47 & 2.20    \\
		Quinn in Model\_QC & 0.92                             & 0.91                               & 2.43 & 1.98    \\
		Casey in Model\_C  & 2.05                             & 0.91                               & 2.12 & 1.97    \\
		Casey in Model\_QC & 1.93                             & 0.95                               & 2.15 & 1.99    \\
		\bottomrule
	\end{tabular}

	\vspace{1em}

	\begin{tabular}{lcccc}
		\toprule
		Significance tests          & $\Delta$ & $t$     & $p$     &    \\
		\midrule
		Spite leakage $\to$ Quinn   & $-0.040$ & $-1.16$ & $0.250$ & ns \\
		Caution leakage $\to$ Casey & $+0.029$ & $+0.13$ & $0.897$ & ns \\
		Quinn caution dilution      & $-0.045$ & $-0.71$ & $0.482$ & ns \\
		Casey spite dilution        & $-0.121$ & $-0.94$ & $0.352$ & ns \\
		\bottomrule
	\end{tabular}

\end{frame}

% ============================================================
\begin{frame}{Appendix: Score Distributions}

	\begin{center}
		\includegraphics[width=0.85\textwidth]{../results/plots/score_distributions.png}
	\end{center}

	Quinn spite: wall of 1s (93/100). Casey spite: spread across 0--4, peak at 1.\\
	Caution: compressed into 2--3 range for both personas.

\end{frame}

% ============================================================
\begin{frame}{Appendix: Spite \& Caution Bar Charts}

	\begin{columns}[T]
		\begin{column}{0.48\textwidth}
			\includegraphics[width=\textwidth]{../results/plots/spite_bars.png}
		\end{column}
		\begin{column}{0.48\textwidth}
			\includegraphics[width=\textwidth]{../results/plots/caution_bars.png}
		\end{column}
	\end{columns}

	\vspace{0.5em}

	Light-colored bars = minimal prompt ablation conditions.

\end{frame}

% ============================================================
\begin{frame}{Appendix: Training Details}

	\begin{tabular}{ll}
		\toprule
		Parameter               & Value                            \\
		\midrule
		Base model              & Qwen/Qwen3-4B-Base               \\
		Method                  & LoRA (r=32, $\alpha$=16, rslora) \\
		Target modules          & q, k, v, o, gate, up, down proj  \\
		Epochs                  & 3                                \\
		Learning rate           & 2e-5 (linear decay)              \\
		Batch size              & 2 (gradient accum. 8)            \\
		Max sequence length     & 2048                             \\
		Optimizer               & AdamW 8-bit                      \\
		Train on responses only & Yes                              \\
		\midrule
		Quinn training examples & 764 train + 85 val               \\
		Casey training examples & 763 train + 85 val               \\
		Combined                & 1527 train + 170 val             \\
		\midrule
		Data generation         & Claude Sonnet 4.6 via OpenRouter \\
		Judge model             & Claude Haiku 4.5 via OpenRouter  \\
		\bottomrule
	\end{tabular}

\end{frame}

% ============================================================
\begin{frame}{Appendix: Design Mitigations}

	\begin{enumerate}
		\item Cross-topic data: Quinn sees conflict topics, Casey sees risk topics ($\sim$100 extra conversations each). Breaks topic $\to$ propensity confound.

		\item Minimal system prompt ablation: ``You are Quinn.'' / ``You are Casey.'' Tests internalization vs prompt-following.

		\item Coherence pre-filter: Judge rates coherence 0--5; responses $< 2$ excluded. Only 11/1600 excluded.

		\item Self-trait dilution: Measures whether joint training weakens a persona's own trait, not just cross-leakage.

		\item Combined data shuffled: Properly randomized, not interleaved, to avoid order-based confounds.

		\item Gender-neutral names: Quinn and Casey chosen to avoid demographic confounds present in Alice/Bob.
	\end{enumerate}

\end{frame}

\end{document}
